% Edited from https://latextemplates.com/template/journal-article
% by Vel (vel@latextemplates.com)
%
% CC BY-NC-SA 4.0 (https://creativecommons.org/licenses/by-nc-sa/4.0/)

\documentclass[letterpaper, 10pt, twoside]{LTJournalArticle}
\addbibresource{references.bib}
\runninghead{FINDS: The Fish INteraction and Dynamics Simulator}

\usepackage{times}

\usepackage{lipsum}
\setcounter{secnumdepth}{0}
\setcounter{page}{1}

\title{FINDS: An N-Body Simulator for Computing the Evolution of Large Schools}

\author{
  Mufaro Machaya\textsuperscript{1},
  Mohamed Niged Mabrouk\textsuperscript{2}, and
  Daniel Floryan\textsuperscript{2}
  \thanks{Corresponding author:
    \href{mailto:dfloryan@uh.edu}{dfloryan@uh.edu}} \\
}

\date{
  August 2026 \vspace{1em} \\
  \footnotesize\textsuperscript{1}University of British Columbia \\
  \footnotesize\textsuperscript{2}University of Houston
}

\renewcommand{\maketitlehookd}{
  \begin{abstract}
    \noindent
    Animals travelling in cohesive groups yields many collective and individual
    benefits, such as security against predators and energy efficiency. The
    replication of such benefits for analogous, biology-inspired technologies
    like autonomous drone systems requires keen understanding of the cohesion
    and evolution of systems of these groups, but these investigations are
    fundamentally limited to small sizes due to the high computational
    complexity in computing pairwise dynamics at order $N^2$.
    As real schools of fish have been recorded to contain millions of members,
    FINDS, or Fish INteraction and Dynamics Simulator, was created to
    efficiently simulate the dynamics of very large schools of fish using
    common computational techniques employed in $N$-body simulation. FINDS is
    implemented as a containerized, modular software suite written
    predominantly in C (with Python for post-processing) that handles three
    core simulation processes in order. First, FINDS contains routines for
    computing the time-derivative of a school, which is a function of the sum
    of the pairwise interactions for each member of the school. As such, FINDS
    contains implementations of Brute Force computation, the Barnes-Hut
    algorithm, and the Fast Multipole Method. Second, FINDS contains custom
    routines for numerically integrating these systems using embedded
    Runge-Kutta schemes for adaptive time-stepping. Third, FINDS contains a
    small library of post-processing modules available to end users to
    stream-line the analysis of their generated data.
  \end{abstract}
}

\begin{document}

\maketitle

\section{Introduction}

\lipsum[1]

\section{Mathematical Model of Swimmers}

FINDS is developed around the far-field source/sink dipole model designed by
Mabrouk and Floryan for modelling how passive hydrodynamic interactions at far
distances effects the evolution of a greater school \cite{mabrouk2024,mabrouk2025}.
Under this model, each fish $i$ is defined as a dipole consisting of a fluid
source at the head of the fish $\mathbf{x}_{f,i}$, and a fluid sink at the
tail $\mathbf{x}_{b,i}$, separated by the length of the fish $\ell_i$. The
source and the sink of each fish are defined to have the same volumetric flow
rate, $\sigma_i$, to ensure the conservation of fluid by making the sources and
sinks identical up to a point.

\section{Literature Review and Similar Work}

\lipsum[1]

\section{Numerical Methods}

\lipsum[1]

\section{Software Architecture}

\lipsum[1]

\section{Validation}

\lipsum[1]

\section{Performance Evaluation}

\lipsum[1]

\section{Results}

\lipsum[1]

\section{Discussion}

\lipsum[1]

\section{Limitations}

\lipsum[1]

\section{Future Work}

\lipsum[1]

\section{Acknowledgements}

In the development of FINDS, Alvin Ng's \texttt{grav\_sim} was used extensively
as a reference for implementing a fast $N$-body simulator in C, particularly
in the implementation of the Barnes-Hut algorithm. In addition, the direct
advice of Alvin Ng was instrumental to the implementation of linear octrees
and parallelization.

Furthermore, generative artificial intelligence (OpenAI's ChatGPT-5.5 and
Google's Gemini 3.3) were used heavily for initial code generation before
being thoroughly tested and validated to ensure correctness. Artificial
intelligence was not used in developing the content of this report, including
the abstract, sections, and figures.

\section{Bibliography}

\printbibliography

\section*{Appendix}

\end{document}
